\chapter{Literature Review}
\label{chap:lr}
\chaptermark{Literature Review}

This chapter reviews the existing work this project uses. The work covered the security of the system as a whole: signing in, what each user may do, how their actions are recorded, and the security of the software the system runs on. A systematic review of access control in electronic health record systems groups this field into identification, authentication, authorisation and accountability \cite{cobrado2024}, and those four categories match the strands of work described in Chapter \ref{chap:met}.

The review is weighted towards authorisation, because that is what most of this project's findings are about. Section \ref{sec:lr-rbac-abac} covers the models of authorisation and what each can express. Section \ref{sec:lr-standards} covers the standards that apply to health information systems, which between them address identification, authentication, transport and audit. Section \ref{sec:lr-threat} covers threat modelling. Section \ref{sec:lr-verification} covers how access control is verified, which is where the main result of this project lies.

\section{Access Control Models: RBAC and ABAC}
\label{sec:lr-rbac-abac}

Access control decides which users may do what, and to which records. Two models are treated as the classic approaches, and FHIR's security guidance describes both for health systems \cite{hl7fhir}. In role-based access control (RBAC), permissions are grouped into roles, each role describes a job function, and users are given roles. If a user's role holds the permission, access is granted \cite{hl7fhir}. Permissions follow the job rather than the person, which suits a hospital, where roles such as surgeon, anaesthetist and finance officer already exist before any software is written.

In attribute-based access control (ABAC), the decision is made from attributes and conditions rather than from roles alone. The National Institute of Standards and Technology (NIST) describes ABAC as a method in which authorisation is determined by evaluating attributes of the subject, attributes of the object, the requested operation, and in some cases the conditions of the environment, against rules that describe which operations a given set of attributes allows \cite{nist80162}. NIST places ABAC on a spectrum of logical access control that also includes access control lists and role-based access control, so it is a development of earlier approaches rather than a replacement for them \cite{nist80162}. The practical gain is that rules may refer to the record itself. ``A surgeon may create a case'' is a role rule. ``A surgeon may edit this case'' is not, because the answer depends on the case, who it is assigned to, not on the surgeon.

Both standards recognise that a decision may also depend on the circumstances of the request, and both make room for it in their models. NIST calls this environment conditions and lists it as one of the four inputs to a decision, alongside rules, subject attributes and object attributes \cite{nist80162}. FHIR uses the term transaction context and names, among its examples, the workflow state of the request \cite{hl7fhir}. The workflow position of a case is therefore recognised in both, though under different names and in different parts of the model.

Work on access control in health systems has followed the same direction. A systematic review of twenty studies on electronic health record systems found attribute-based access control to be the most common authorisation mechanism, and grouped the mechanisms proposed into identification, authentication, authorisation and accountability \cite{cobrado2024}. The review also names what the literature has not resolved: multi-factor authentication, emergency access, patient consent and accountability are identified as gaps, with accountability addressed by only two of the twenty studies. Others have combined the two models directly. Atlam and Yang propose a healthcare model integrating role-based categorisation with attribute-based control over user attributes and environmental factors, observing that models which allocate permissions by predetermined roles or policies adapt poorly to the changing circumstances of a working healthcare system \cite{atlam2025}.

The three inputs are not alternatives. Roles carry the decisions that follow the job, attributes carry those that depend on the individual record, and the case's position in the workflow carries those that depend on how far it has progressed. What neither standard describes is how they combine in practice, what it costs to enforce them consistently, or how such rules can be checked once built. Chapters \ref{chap:impl} and \ref{chap:result} report that.

\section{Security Standards for Health Information Systems}
\label{sec:lr-standards}

Health information systems carry data that is sensitive in ways ordinary business systems are not. FHIR's security guidance was the health-specific source used here. It is unusual in addressing implementation directly: it recommends OAuth for authenticating users, states that production data should be sent over TLS, describes how audit events should be recorded, and sets out how a server should respond when access is refused \cite{hl7fhir}.

No Cambodian standard exists for health information security that a hospital can adopt. The nearest domestic guidance is the National Bank of Cambodia's Technology and Cyber Risk Management Guidelines, issued in January 2026 \cite{nbc2026}. These apply to banks rather than hospitals, so this system is not subject to them. They are, however, the only Cambodian document setting out detailed security requirements for systems holding sensitive data, and the requirements used here concern the protection of that data rather than financial services.

The guidelines require role-based access control with least privilege, a maintained access control matrix, segregation of duties so that no single person can enter, authorise and complete a transaction, audit logs that record both successful and refused access and cannot be altered, encryption of data in transit and at rest, and vulnerability assessment using both automated tools and manual techniques \cite{nbc2026}. They also note that Cambodia's adoption of technology is outpacing its regulatory frameworks, and direct institutions to adopt recognised standards where no local regulation exists \cite{nbc2026}.

This project follows that direction. Two general frameworks were used alongside FHIR. The NIST Cybersecurity Framework organises security work into six functions, Govern, Identify, Protect, Detect, Respond and Recover, and was used to structure the phases of the project rather than as a control set implemented in full \cite{nistcsf2}. The Open Worldwide Application Security Project (OWASP) publishes a list of the ten most common categories of web application security risk, used here as the checklist for reviewing the code \cite{owasp2021}. Each source serves one purpose: FHIR gives health-specific implementation guidance, NIST SP 800-162 the access control model, the National Bank of Cambodia guidelines a domestic benchmark, the NIST Cybersecurity Framework a structure for the work, and the OWASP Top 10 a checklist for the code.

None of these covers the part that matters most for a workflow system. They describe how to authenticate users, how to protect data in transit, what kinds of flaws to look for, and how to organise the work. None states who in a particular hospital may perform a particular action on a particular case, because that depends on how that hospital divides responsibility. Those rules must be drawn from the workflow itself, and the standards give no method for doing so, nor for checking the result.

\section{Threat Modelling Approaches}
\label{sec:lr-threat}

Threat modelling is the practice of working out what could go wrong with a system, without testing it. It asks what someone might attempt against each part of a system, and what would follow if they succeeded. OWASP treats it as a design activity. Its category for insecure design calls for greater use of threat modelling in authentication, access control and business logic, and separates design flaws from implementation defects, on the grounds that a design flaw cannot be fixed by careful coding \cite{owasp2021}.

STRIDE is the most widely used method. It groups threats into six kinds \cite{stride}: spoofing, where someone acts as another user; tampering, where data is changed without permission; repudiation, where a user denies an action and no record contradicts them; information disclosure, where data reaches someone who should not see it; denial of service, where the system is made unavailable; and elevation of privilege, where a user gains rights they were not granted. Each part of a system is examined against all six.

The value of the method is coverage, not discovery. Working through six categories for every component makes it harder to miss a whole class of threat. What STRIDE does not do is say which threats matter most, or supply the rules that would prevent them. It can tell you a system may be vulnerable to tampering. It cannot tell you which staff member should lose the ability to edit a record once another department has reviewed it. That judgement comes from the workflow, not from the method.

Threat modelling is usually described as a design-time activity, carried out before a system is built. It can also be applied to a system that already exists. In that case it organises and classifies what is known, rather than predicting what has not yet been found. Chapter \ref{chap:met} describes which of these applies to this project.

\section{Automated and Manual Security Verification}
\label{sec:lr-verification}

Writing down who may do what is one task. Checking that the software actually does it is another. A rule table states what should happen, but it says nothing about what the code does. The two can disagree without anyone noticing, because when a system allows something it should have blocked, nothing appears to go wrong. Not every strand of this work presents the same difficulty. Dependency versions, container images and secrets left in a repository can each be checked by tools, because each has a recognisable form. Access control does not, which is why it is the strand this section concentrates on.

Automated tools are the usual first step, and their main difficulty is well documented. Sun, Xu and Su point out that access control rules are specific to each application \cite{sun2011}. Unlike an injection flaw, which looks much the same in any program, there is no general pattern for a tool to search for. The rules have to be supplied, and writing them by hand is slow and tends to produce rules that are missing, incomplete or wrong. Their own tool works around this by inferring the rules from the source code: it builds a list of the pages each role can open, compares those lists to find pages only some roles should reach, and then tests whether other users can reach them anyway.

Inferring the rules from code is where the method reaches its limit. Zhu and colleagues tested six open source applications and found that the assumptions built into earlier tools held those tools back \cite{zhu2015}. They tried a different approach, asking developers to supply the rules and turning those answers into patterns the tool could then apply. This found twenty vulnerabilities the automatic methods had missed. A survey published in 2024 reaches the same conclusion, comparing the methods available and listing the problems that remain unsolved \cite{anas2024}. The pattern across all three is that access control rules come from the organisation rather than from the code, and that a tool performs better once a person supplies them.

This is why automated analysis cannot be the whole of the work, and the guidance says as much. OWASP asks developers to write tests for access control \cite{owasp2021}. The National Bank of Cambodia's guidelines require both automated tools and manual techniques \cite{nbc2026}. Neither explains how to test permissions that change as a case moves from one stage to the next. Chapter \ref{chap:result} reports what each method found when both were used on the same system in the same period.
